\documentclass[11pt,oneside]{article}

\pdfminorversion=7
\usepackage{QuizStyle}

\hypersetup{
    pdftitle={20260922 Quiz for MATH301 Modern Algebra I},
    pdfauthor={Jungwoo Kim},
    pdfsubject={Quiz for MATH301 Modern Algebra I},
    pdfcreator={LaTeX}
}

\begin{document}

\quiztitle{20260922 Quiz}

\begin{exercise}{1.1.21}
Let $G$ be a finite group and let $x$ be an element of $G$ of order $n$. Prove that if $n$ is odd, then
\[
    x=(x^2)^k
\]
for some $k$.
\end{exercise}

\begin{solution}
Since $n$ is odd, $k=(n+1)/2$ is an integer. As $|x|=n$,
\[
    (x^2)^k=x^{n+1}=x^n x=x.
\]
\end{solution}

\begin{exercise}{1.1.25}
Prove that if $x^2=1$ for all $x\in G$, then $G$ is abelian.
\end{exercise}

\begin{solution}
For every $x\in G$, the hypothesis gives $x^{-1}=x$. Hence for arbitrary $a,b\in G$,
\[
    ab=(ab)^{-1}=b^{-1}a^{-1}=ba.
\]
Thus $G$ is abelian.
\end{solution}

\begin{exercise}{1.1.31}
Prove that any finite group $G$ of even order contains an element of order $2$.

\textit{Hint.} Let
\[
    t(G)=\{g\in G\mid g\neq g^{-1}\}.
\]
Show that $t(G)$ has an even number of elements and every nonidentity element of $G\setminus t(G)$ has order $2$.
\end{exercise}

\begin{solution}
If $g\in t(G)$, then $g^{-1}\in t(G)$ and $g\neq g^{-1}$. Thus the elements of $t(G)$ occur in disjoint pairs $\{g,g^{-1}\}$, so $|t(G)|$ is even. If $x\in G\setminus t(G)$, then $x=x^{-1}$; hence every nonidentity element of $G\setminus t(G)$ satisfies $x^2=1$ and has order $2$. Since $|G|$ and $|t(G)|$ are even, $|G\setminus t(G)|$ is even. As $1\in G\setminus t(G)$, this set contains at least one nonidentity element, and therefore $G$ contains an element of order $2$.
\end{solution}

\begin{exercise}{1.1.34}
If $x$ is an element of infinite order in $G$, prove that the elements $x^n$, $n\in\Z$, are all distinct.
\end{exercise}

\begin{solution}
Suppose $x^m=x^n$ for some $m,n\in\Z$. If $m\neq n$, assume without loss of generality that $m>n$. Then
\[
    x^{m-n}=1,
\]
with $m-n>0$, contradicting that $x$ has infinite order. Hence $m=n$, so the powers $x^n$ are all distinct.
\end{solution}

\begin{exercise}{1.2.3}
Use the presentation
\[
    D_n=\langle r,s\mid r^n=s^2=1,\ rs=sr^{-1}\rangle
\]
to show that every element of $D_n$ which is not a power of $r$ has order $2$. Deduce that $D_n$ is generated by the two elements $s$ and $sr$, both of which have order $2$.
\end{exercise}

\begin{notation}
Following MATH000, $D_n$ denotes the dihedral group of symmetries of a regular $n$-gon, so $|D_n|=2n$. Dummit--Foote denotes the same group by $D_{2n}$.
\end{notation}

\begin{solution}
Every element of $D_n$ which is not a power of $r$ has the form $sr^i$. From $rs=sr^{-1}$ we obtain $r^i s=sr^{-i}$, so
\[
    (sr^i)^2=sr^isr^i=s(sr^{-i})r^i=1.
\]
Since $sr^i\neq1$, it has order $2$. In particular, $s$ and $sr$ have order $2$. Also
\[
    s(sr)=r,
\]
so $s$ and $sr$ generate both $r$ and $s$, hence generate $D_n$.
\end{solution}

\begin{exercise}{1.2.6}
Let $x$ and $y$ be elements of order $2$ in any group $G$. Prove that if $t=xy$, then
\[
    tx=xt^{-1}.
\]
Thus, if $n=|xy|<\infty$, then $x$ and $t$ satisfy the same relations in $G$ as $s$ and $r$ do in $D_n$.
\end{exercise}

\begin{solution}
Since $x^{-1}=x$ and $y^{-1}=y$,
\[
    t^{-1}=(xy)^{-1}=yx.
\]
Hence
\[
    tx=(xy)x=xyx=x(yx)=xt^{-1}.
\]
If $n=|t|<\infty$, then
\[
    t^n=x^2=1,
    \qquad
    tx=xt^{-1},
\]
which are the defining relations for $r$ and $s$ in $D_n$ under $r\leftrightarrow t$ and $s\leftrightarrow x$.
\end{solution}

\begin{exercise}{1.2.7}
Show that
\[
    \langle a,b\mid a^2=b^2=(ab)^n=1\rangle
\]
gives a presentation for $D_n$ in terms of the two generators $a=s$ and $b=sr$ of order $2$.

\textit{Hint.} Show that the relations for $r$ and $s$ follow from the relations for $a$ and $b$, and conversely, that the relations for $a$ and $b$ follow from those for $r$ and $s$.
\end{exercise}

\begin{solution}
Set $a=s$ and $b=sr$. Then
\[
    a^2=s^2=1,
    \qquad
    b^2=(sr)^2=1,
    \qquad
    ab=s(sr)=r,
\]
so $(ab)^n=r^n=1$.

Conversely, assume $a^2=b^2=(ab)^n=1$ and set
\[
    s=a,
    \qquad
    r=ab.
\]
Then $s^2=1$ and $r^n=1$. Since $a^{-1}=a$ and $b^{-1}=b$,
\[
    r^{-1}=(ab)^{-1}=ba,
\]
and therefore
\[
    rs=(ab)a=aba=a(ba)=sr^{-1}.
\]
Thus the two sets of generators and relations determine one another, giving the stated presentation of $D_n$.
\end{solution}

\begin{exercise}{1.2.15}
Find a set of generators and relations for $\Z/n\Z$.
\end{exercise}

\begin{solution}
A generating set is $\{[1]_n\}$, and the defining relation is
\[
    n[1]_n=[0]_n.
\]
Equivalently, in multiplicative presentation notation,
\[
    \Z/n\Z\cong\langle x\mid x^n=1\rangle.
\]
\end{solution}

\begin{exercise}{1.3.5}
Find the order of
\[
    (1\ 12\ 8\ 10\ 4)(2\ 13)(5\ 11\ 7)(6\ 9).
\]
\end{exercise}

\begin{solution}
The displayed cycles are disjoint and have lengths $5,2,3,2$. Hence the order is
\[
    \lcmop(5,2,3,2)=30.
\]
\end{solution}

\begin{exercise}{1.3.13}
Show that an element has order $2$ in $S_n$ if and only if its cycle decomposition is a product of commuting $2$-cycles.
\end{exercise}

\begin{solution}
Let $\sigma\in S_n$ have order $2$. Then $\sigma^2=1$. If a nontrivial cycle in the disjoint cycle decomposition of $\sigma$ had length $m\ge3$, applying $\sigma$ twice would not fix the elements in that cycle. Hence every nontrivial cycle has length $2$. Since disjoint cycles commute, $\sigma$ is a product of commuting $2$-cycles.

Conversely, if $\sigma$ is a nonempty product of disjoint $2$-cycles, then the cycles commute and each has square $1$, so $\sigma^2=1$. Since $\sigma\neq1$, it follows that $|\sigma|=2$.
\end{solution}

\begin{exercise}{1.3.14}
Let $p$ be a prime. Show that an element has order $p$ in $S_n$ if and only if its cycle decomposition is a product of commuting $p$-cycles. Show by an explicit example that this need not be the case if $p$ is not prime.
\end{exercise}

\begin{solution}
Suppose $|\sigma|=p$. Then $\sigma^p=1$. If a nontrivial cycle in the disjoint cycle decomposition of $\sigma$ has length $m$, then an $m$-cycle returns to the identity after $p$ applications only when $m\mid p$. Since $p$ is prime and $m>1$, we have $m=p$. Thus $\sigma$ is a product of disjoint, hence commuting, $p$-cycles.

Conversely, a nonempty product of disjoint $p$-cycles has $p$th power equal to $1$, while no positive power smaller than $p$ is the identity. Hence its order is $p$.

The primality assumption is necessary. For example,
\[
    \sigma=(1\ 2)(3\ 4\ 5\ 6)
\]
has order $4$, but its cycle decomposition is not a product solely of $4$-cycles.
\end{solution}

\begin{exercise}{1.3.15}
Prove that the order of an element in $S_n$ equals the least common multiple of the lengths of the cycles in its cycle decomposition.

\textit{Hint.} Use Exercise 1.3.10 and Exercise 1.1.24.
\end{exercise}

\begin{solution}
Write the disjoint cycle decomposition as
\[
    \sigma=\sigma_1\sigma_2\cdots\sigma_r,
\]
where $\sigma_i$ has length $m_i$. Since the cycles are disjoint, they commute, and therefore
\[
    \sigma^k=\sigma_1^k\sigma_2^k\cdots\sigma_r^k.
\]
This is the identity if and only if each $\sigma_i^k$ is the identity, which holds if and only if $m_i\mid k$ for every $i$. Hence the least positive such $k$ is
\[
    \lcmop(m_1,\ldots,m_r).
\]
\end{solution}

\begin{exercise}{1.3.18}
Find all numbers $n$ such that $S_5$ contains an element of order $n$.

\textit{Hint.} Use Exercise 1.3.15.
\end{exercise}

\begin{solution}
The possible nontrivial disjoint cycle types in $S_5$ have lengths
\[
    2,\quad 2+2,\quad 3,\quad 3+2,\quad 4,\quad 5.
\]
By Exercise 15 their orders are respectively
\[
    2,\quad 2,\quad 3,\quad 6,\quad 4,\quad 5.
\]
Including the identity, the possible orders are therefore
\[
    \boxed{1,2,3,4,5,6}.
\]
\end{solution}

\begin{exercise}{1.6.4}
Prove that the multiplicative groups $\R\setminus\{0\}$ and $\C\setminus\{0\}$ are not isomorphic.
\end{exercise}

\begin{solution}
The element $i\in\C\setminus\{0\}$ has order $4$. If the two groups were isomorphic, $\R\setminus\{0\}$ would also contain an element of order $4$, since isomorphisms preserve element orders. But if $x\in\R\setminus\{0\}$ satisfies $x^4=1$, then $x=\pm1$, so $x$ has order $1$ or $2$. This is a contradiction.
\end{solution}

\begin{exercise}{1.6.14}
Let $G$ and $H$ be groups and let $\varphi\colon G\to H$ be a homomorphism. Define the \emph{kernel} of $\varphi$ to be
\[
    \ker\varphi=\{g\in G\mid \varphi(g)=1_H\}.
\]
Prove that $\ker\varphi$ is a subgroup of $G$ (cf.\ Exercise 1.1.26). Prove that $\varphi$ is injective if and only if $\ker\varphi$ is the identity subgroup of $G$.
\end{exercise}

\begin{solution}
Since $\varphi(1_G)=1_H$, we have $1_G\in\ker\varphi$. If $a,b\in\ker\varphi$, then
\[
    \varphi(ab)=\varphi(a)\varphi(b)=1_H,
\]
and
\[
    \varphi(a^{-1})=\varphi(a)^{-1}=1_H.
\]
Thus $\ker\varphi$ is closed under products and inverses, so $\ker\varphi\le G$.

If $\varphi$ is injective and $g\in\ker\varphi$, then
\[
    \varphi(g)=1_H=\varphi(1_G),
\]
so $g=1_G$. Hence $\ker\varphi=\{1_G\}$. Conversely, suppose $\ker\varphi=\{1_G\}$ and $\varphi(a)=\varphi(b)$. Then
\[
    \varphi(ab^{-1})=\varphi(a)\varphi(b)^{-1}=1_H,
\]
so $ab^{-1}=1_G$, and therefore $a=b$. Hence $\varphi$ is injective.
\end{solution}

\begin{exercise}{1.6.19}
Let
\[
    G=\{z\in\C\mid z^n=1\text{ for some }n\in\N\}.
\]
Prove that for any fixed integer $k>1$, the map from $G$ to itself defined by
\[
    z\mapsto z^k
\]
is a surjective homomorphism but is not an isomorphism.
\end{exercise}

\begin{solution}
Define $\varphi(z)=z^k$. Since multiplication in $\C$ is commutative,
\[
    \varphi(zw)=(zw)^k=z^kw^k=\varphi(z)\varphi(w),
\]
so $\varphi$ is a homomorphism.

Let $w\in G$. Choose $z\in\C$ with $z^k=w$. If $w^n=1$, then
\[
    z^{kn}=(z^k)^n=w^n=1,
\]
so $z\in G$ and $\varphi(z)=w$. Thus $\varphi$ is surjective.

Finally, let $\zeta=e^{2\pi i/k}$. Then $\zeta\neq1$, $\zeta\in G$, and
\[
    \varphi(\zeta)=\zeta^k=1=\varphi(1).
\]
Hence $\varphi$ is not injective and therefore is not an isomorphism.
\end{solution}

\begin{exercise}{1.6.22}
Let $A$ be an abelian group and fix some $k\in\Z$. Prove that the map
\[
    a\mapsto a^k
\]
is a homomorphism from $A$ to itself. If $k=-1$, prove that this homomorphism is an isomorphism, i.e., an automorphism of $A$.
\end{exercise}

\begin{solution}
For $a,b\in A$, the elements commute, so the exponent laws give
\[
    (ab)^k=a^kb^k.
\]
Hence the map $\varphi(a)=a^k$ satisfies
\[
    \varphi(ab)=\varphi(a)\varphi(b),
\]
and is a homomorphism.

If $k=-1$, then $\varphi(a)=a^{-1}$ and
\[
    \varphi(\varphi(a))=(a^{-1})^{-1}=a.
\]
Thus $\varphi^{-1}=\varphi$, so $\varphi$ is an automorphism of $A$.
\end{solution}

\begin{exercise}{1.6.25}
Let $n\in\N$, let $r$ and $s$ be the usual generators of $D_n$, and let $\theta=2\pi/n$.
\begin{exercisesubparts}
    \item Prove that
    \[
        \begin{pmatrix}
            \cos\theta & -\sin\theta\\
            \sin\theta & \cos\theta
        \end{pmatrix}
    \]
    is the matrix of the linear transformation which rotates the $xy$-plane about the origin in a counterclockwise direction by $\theta$ radians.

    \item Prove that the map $\varphi\colon D_n\to\GL_2(\R)$ defined on generators by
    \[
        \varphi(r)=
        \begin{pmatrix}
            \cos\theta & -\sin\theta\\
            \sin\theta & \cos\theta
        \end{pmatrix},
        \qquad
        \varphi(s)=
        \begin{pmatrix}
            0&1\\
            1&0
        \end{pmatrix}
    \]
    extends to a homomorphism of $D_n$ into $\GL_2(\R)$.

    \item Prove that the homomorphism $\varphi$ in part (b) is injective.
\end{exercisesubparts}
\end{exercise}

\begin{solution}
Denote the matrix in part (a) by $\mathbf R$ and the matrix assigned to $s$ in part (b) by $\mathbf S$.
\begin{exercisesubparts}
    \item A counterclockwise rotation by $\theta$ sends
    \[
        \begin{pmatrix}1\\0\end{pmatrix}
        \mapsto
        \begin{pmatrix}\cos\theta\\\sin\theta\end{pmatrix},
        \qquad
        \begin{pmatrix}0\\1\end{pmatrix}
        \mapsto
        \begin{pmatrix}-\sin\theta\\\cos\theta\end{pmatrix}.
    \]
    These images are the columns of $\mathbf R$, proving part (a).

    \item Since $\mathbf R$ is rotation by $2\pi/n$, $\mathbf R^n=\mathbf I$, and directly $\mathbf S^2=\mathbf I$. Moreover,
    \[
        \mathbf R\mathbf S
        =\begin{pmatrix}-\sin\theta&\cos\theta\\\cos\theta&\sin\theta\end{pmatrix}
        =\mathbf S\mathbf R^{-1}.
    \]
    Thus $\mathbf R$ and $\mathbf S$ satisfy $r^n=s^2=1$ and $rs=sr^{-1}$. Hence $r\mapsto\mathbf R$, $s\mapsto\mathbf S$ extends to a homomorphism $\varphi\colon D_n\to\GL_2(\R)$.

    \item By Exercise 1.6.14, it suffices to show $\ker\varphi=\{1\}$. Every element of $D_n$ is $r^i$ or $sr^i$ with $0\le i<n$. If $\mathbf R^i=\mathbf I$, then rotation by $2\pi i/n$ is the identity, so $i=0$. Also
    \[
        \det(\mathbf S\mathbf R^i)=\det(\mathbf S)\det(\mathbf R^i)=-1\neq\det(\mathbf I),
    \]
    so $\mathbf S\mathbf R^i\neq\mathbf I$ for every $i$. Hence $\ker\varphi=\{1\}$, and $\varphi$ is injective.
\end{exercisesubparts}
\end{solution}

\end{document}
